\subsection{معرفی و انگیزه}
این ماژول در چارچوب پژوهش حاضر پیاده‌سازی و ارزیابی شده و فاقد مقاله مستقل جداگانه در ادبیات تحقیق است.

کاهش ابعاد ماتریس‌های وزن در مدل‌های ترنسفورمر از طریق تجزیه کم‌مرتبه (\lr{Low-Rank Decomposition}) یکی از روش‌های کلیدی جهت فشرده‌سازی و کاهش پیچیدگی محاسباتی است. با این وجود، روش‌های سنتی مانند تجزیه مقادیر منفرد قطعی (\lr{Deterministic SVD}) بدون آگاهی از داده‌های ورودی، وزن‌های تمام ابعاد را با اهمیت یکسان در نظر می‌گیرند و منجر به تخریب بازنمایی ویژگی‌های با انرژی بالا در فعال‌سازی‌ها می‌شوند.

ماژول \lr{optimal\_subspace\_projection} با نام فارسی «پروجکشن زیرفضای بهینه» به عنوان یکی از اجزای پیشرفته موتور فشرده‌سازی (\lr{Compression Engine})، رویکرد پروجکشن زیرفضای آگاه از فعال‌سازی (\lr{Activation-Aware Subspace Projection}) را از طریق ترکیبی از فاکتورگیری چولسکی (\lr{Cholesky Preconditioning}) و تجزیه مقادیر منفرد تصادفی (\lr{Randomized SVD with Power Iterations}) پیاده‌سازی می‌کند. این روش تضمین می‌کند که زیرفضای پروجکشن به دست آمده، واریانس بیشینه پاسخ‌های فعال‌سازی را پوشش داده و بازسازی ماتریس وزن را با حداقل خطای فروبنیوس همراه سازد.
